\documentclass[10pt, a4paper]{article}

\usepackage[margin=1.5cm]{geometry}
\usepackage[utf8]{inputenc}
\usepackage[spanish,es-noshorthands]{babel}
\usepackage{amssymb, amsmath, amsbsy}
\usepackage{verbatim}
\usepackage{mathpazo}
\usepackage{tasks}
\usepackage{enumerate}
\usepackage[pdftex]{hyperref}
\hypersetup{pdftitle={Examen de Investigación Operativa },pdfauthor={Luis Carrillo Gutiérrez}}

\fontfamily{cmss}\selectfont
\renewcommand{\familydefault}{\sfdefault}
\pagestyle{empty}

\begin{document}
\begin{center}
\subsection*{Tercer Examen de Investigación Operativa} 
\end{center}

\noindent Nombres / Apellidos : \rule{10cm}{0.4pt} \\ 
Fecha : \rule{2cm}{0.4pt}  / \rule{2cm}{0.4pt} / \rule{2cm}{0.4pt} \hspace{3.4cm} Duración: {\small 90 minutos} \\

\begin{enumerate}
\item{Proponer una definición para los siguientes conceptos :
\begin{itemize}
\item{Casillero básico} 
\item{Arco básico}
\item{Ciclo marginal}
\item{Costo marginal}
\end{itemize}
}
\item{Hallar el valor de las \emph{variables duales} del siguiente problema e intérpretar adecuadamente los resultados :
$$Max\;\mathbf{Z} = 200\,x_{1} + 220\,x_{2} + 240\,x_{3}$$
s.a
\begin{eqnarray*}
2 x_{1} + x_{2} + 2 x_{3} & \leq & 4800 \\
x_{1} + x_{2} + 2 x_{3} & \leq & 3400 \\
2 x_{1} + x_{2} + x_{3} & \leq & 4000 \\
x_{1}\;\geq 0,\quad x_{2}\;\geq 0,\; & x_{3} & \geq 0 
\end{eqnarray*}
}
\item{Hallar la solución óptima del siguiente problema de transporte ({\small Los datos de cada casillero son los costos}).

\begin{center}
	\begin{tabular}{|c|c|c|c|c|}
		\hline 
		\empty & \multicolumn{4}{c|}{$b_{i}$} \\
		\hline
		$a_{i}$ & 240 & 200 & 250 & 100 \\
		\hline 
		160 & 5 & 6 & 9 & 8 \\
		180 & 8 & 8 & 4 & 9 \\
		280 & 9 & 10 & 6 & 7 \\
		200 & 7 & 5 & 8 & 10 \\
		160 & 8 & 9 & 10 & 7 \\
		\hline
	\end{tabular}
\end{center}
}
\item{Se desea asignar $6$ obreros a $6$ puestos de trabajo, la tabla muestra los rendimientos de cada obrero en cada puesto. Hallar la \emph{asignación óptima}.

\begin{center}
	\begin{tabular}{| p{1cm} ||c|c|c|c|c|c|}
		\hline
		\empty & A & B & C & D & E & F \\
		\hline
		1 & 90 & 100 & 92 & 102 & 78 & 78 \\
		2 & 76 & 78 & 94 & 80 & 79 & 84 \\
		3 & 85 & 92 & 86 & 92 & 91 & 80 \\
		4 & 79 & 97 & 84 & 95 & 84 & 76 \\
		5 & -- & 87 & 89 & 80 & 85 & 88 \\
		6 & 80 & 90 & 79 & 82 & 84 & 82 \\
		\hline
	\end{tabular}
\end{center}
}
\end{enumerate}
\end{document}